\chapter{Conclusion}
\label{chap:conclusion}

This report has offered a survey of the state of the art methods in pruning Transformers. We've highlighted that while local scoring approaches may perform well for relatively low compression levels, they are limited at high sparsity levels as they do not take into account the interaction between components that were pruned. As the ratio of compression increases, how the sparsity budget is allocated across the layers becomes of utmost importance and local scoring approaches are unable to do so.

From the studies of the current methods, it can be observed that the Transformer pruning task needs to be viewed from a new perspective which is that the pruning should not be modeled as a ranking task in local, but rather as a configuration searching task in global in future work. This view can provide better insights into the dependencies of pruned parts with each other and lead to an efficient pruning strategy which achieves more significant compression ratio.

In summary, from the disadvantages of local scoring methods to prune Transformer, we can clearly see that when designing a pruning method for Transformer, we should not neglect the interaction and configuration of the Transformer globally. We believe that the future works in this area should be based on searching strategies pruning framework which can properly search the huge configuration space of pruning and attain optimum performance in high sparsity.